\documentclass[../main]{subfiles}
\begin{document}
\ifSubfilesClassLoaded{\setcounter{page}{36}}{}
\section{Необходимый труд и условия для коммунистической деятельности}
\label{sec:06_necessitatem_truda}

В коммунистическом обществе жизнь человека не опосредована частной собственностью и обменом товарами, следовательно, она свободна от всего того, что нужно для обеспечения условий существования частной собственности. Это даёт обществу значительные преимущества и сберегает его рабочие силы. В «Эльберфельдских речах» Энгельс приводил ряд примеров, нагляднейшим образом показывающих, сколько человеческого труда может быть сохранено в обществе, где интересы членов не противоположны, а продукты деятельности не превращаются в товары, противостоящие друг другу на рынке.

Труд, необходимый для функционирования рынка, через который распределяются общественные ресурсы, заменяется более эффективным прямым распределением. Это даёт возможность не совершать действий, бессмысленных с точки зрения организации производства в общественном масштабе и не растрачивать на них время жизни людей. Производить упаковку, не нужную с точки зрения простого хранения продуктов, совершать перевозки, связанные с продажами, а не с доставкой конечному потребителю, производить продукты, которые потом окажутся лишними на рынке и не будут проданы, и уж тем более, производить то, что приносит вред человеку, однако приносит огромные прибыли в обществе товарного производства (наркотики и т.п.), - всё это не имеет никакого смысла в обществе, где производство не носит товарный характер.

Уничтожение частного характера присвоения тоже дает колоссальную экономию живого человеческого труда. Это то же самое, что и товарное производство, но рассмотренное со стороны труда, затраченного на доведение до конечного потребления. Примеры быта и транспорта, о которых мы уже писали, прямо-таки кричат о возможности высвобождения большого количества человеческого    труда,   который     сейчас   расходуется    неэффективным индивидуальным способом.

Современный капитализм требует усилий людей, направленных на поддержание возможности производства прибавочной стоимости. Труда такого типа становится всё больше и больше. Производство прибавочной стоимости с необходимостью дополнилось производством потребления и образа жизни, нацеленного на потребление тех или иных товаров. Уже одно это показывает возрастающие (и значительно возросшие по сравнению с временами Маркса и Энгельса) возможности для освобождения людей, которые смогут на началах коллективной собственности заниматься производством того, что нужно обществу для его развития.

Общественное время тратится на поддержание существования тех ограничений, которые накладывает на жизнь общества частная собственность, и устранение её последствий, вредных для общества, в том числе и преступлений: полиция, регулярная армия, чиновники, разного рода охранники частной собственности. Все эти рабочие силы тоже могут быть освобождены с преодолением     частной   собственности  на    средства   производства    и перенаправлены на выполнение труда, в котором общество непосредственно заинтересовано.

Всё это уже само по себе, даже на базе уже существующих технологий, снизит необходимое общественное рабочее время и создаст простор для существования свободного времени для действительной человеческой жизни, то есть собственно для коммунизма.

На движение в направлении коммунизма указывает ещё одна важная тенденция, которая зреет в недрах капитализма и тормозится им только потому, что капитал является господствующим общественным отношением. «Великая историческая сторона капитала заключается в создании этого прибавочного труда, излишнего с точки зрения одной лишь потребительной стоимости, с точки зрения простого поддержания существования рабочего, и историческое назначение капитала будет выполнено тогда, когда, с одной стороны, потребности будут развиты настолько, что сам прибавочный труд, труд за пределами абсолютно необходимого для жизни, станет всеобщей потребностью, проистекающей из самих индивидуальных потребностей людей, и когда, с другой стороны, всеобщее трудолюбие благодаря строгой дисциплине капитала, через которую прошли следующие друг за другом поколения, разовьется как всеобщее достояние нового поколения, - когда, наконец, это всеобщее трудолюбие, благодаря развитию производительных сила труда, постоянно подстегиваемых капиталом, одержимым беспредельной страстью к обогащению и действующим в таких условиях, в которых он только и может реализовать эту страсть, приведет к тому, что, с одной стороны, владение всеобщим богатством и сохранение его будут требовать от всего общества лишь сравнительно незначительного количества рабочего времени и что, с другой стороны, работающее общество будет по-научному относится к процессу своего прогрессирующего воспроизводства, своего воспроизводства во все возрастающем изобилии; - следовательно, тогда, когда прекратится такой труд, при котором человек сам делает то, что он может заставить вещи делать для себя, для человека» - писал Маркс. Мы же добавим, что это было бы завершением всего периода существования и развития общественного разделения труда.

Идея автоматизированной фабрики, появившаяся ещё во времена Маркса, и идея широкого использования возможностей природы при помощи знаний о ней на благо всех людей, охватывает реально существующее тенденции к созданию условий для коммунизма. Коммунизм будет наступать по мере того, как вещи всё больше и больше будут делать работу, которую они могут сделать для человека и без человека, всё меньше и меньше будет несвободного человеческого труда, необходимого для воспроизводства человеческой жизни, и всё больше управление людьми будет превращаться в управление вещами. Каждый человек всё больше и больше времени сможет отдавать свободной игре физических и умственных сил, или, другими словами, на то, чтобы заниматься тем, что ему «по душе», на пользу всего общества. Последнее и составляет сущность коммунистической деятельности.

Историческая роль капитала как существенного общественного отношения для развития производительных сил не может быть выполнена при капитализме. При безраздельном своём господстве капитал постоянно воспроизводит вещный характер труда, то и дело отбрасывая развитие производительных сил назад там, где наметилась тенденция его преодоления. Поэтому наука при капитализме не может окончательно превратиться в непосредственную производительную силу, а при воспроизводстве вещного характера труда воспроизводятся, в том числе, и самые отсталые его формы. Полное уничтожение капитала как общественного отношения – важнейший вопрос переходного периода к коммунистическому обществу. В этот период нужно сознательно доводить движение капитала по передаче чисто материального производства в ведение машин, до выведения человеческой деятельности за пределы вещного характера труда.

При этом не следует забывать о том, что капитал в то же время воспроизводит и противоположную тенденцию, которая ведёт общество назад к капитализму. Вот с нею, хотя она неразрывно связана с первой, социалистическое общество может и должно открыто бороться, чтобы переход к коммунизму вообще мог состояться. Линия фронта между старым капиталистическим и новым коммунистическим объективно проходит там, где есть налицо все материальные предпосылки, для того, чтобы вещный труд раз и навсегда передать вещам, и там, где этого сделать пока еще нельзя, не доведя труд до крайней степени разделения и создания, таким образом, вещественных предпосылок, чтобы передать его машинам. Там, где возможна полная автоматизация чисто материального производства вещей, капитал как общественное отношение и соответствующий характер труда утрачивает свою актуальность и потому объективно может и должен быть уничтожен. И пока это не будет сделано, капитал все равно будет воспроизводиться как существенное общественное отношение.

Пресловутый вопрос антикоммунистически настроенного обывателя «кто будет мыть унитазы, если каждый будет заниматься любимым делом?» лишается всякого смысла. Его пафос обессмысливается не только потому, что существует огромная разница между тем, чтобы иногда отвлечься от любимого дела для выполнения этой неприятной, но необходимой общественной обязанности, и тем, чтобы быть на всю жизнь прикованным к ней, но и потому, что сейчас уже существуют унитазы, которые сами себя моют. И только господство частной собственности и то, что рабочая сила уборщиков, то есть время жизни людей, оказывается в этих условиях дешевле вещей, мешает их повсеместному распространению.

При коммунизме же человек, то есть время и наполненность его жизни, становится наивысшей ценностью, а все вещи имеют значение только постольку, поскольку они обеспечивают действительную человеческую жизнь каждого члена общества. Сам факт того, что живой человек должен время своей жизни тратить на нетворческий труд, который можно передать машинам, невыносим для общества, где человек является наивысшей ценностью. Такое положение вещей должно быть уничтожено, и только после этого можно будет вести речь о коммунизме, развитом на своём собственном основании.

Однако, следует признать, что в переходный период к коммунизму какое-то время необходимый лишённый полной свободы труд будет сохраняться, пока для его упразднения не разовьются соответствующие технологии и они не будут внедрены повсеместно. Но устранение господства капитала позволит ограничить рабочий день необходимым трудом, исключив труд, имеющий смысл только в рамках движения стоимости к самовозрастанию, результаты которого то и дело уничтожаются во время кризисов и просто в качестве «издержек». На первых порах необходимый (в смысле не свободный ещё в полной мере) труд расширит свои рамки. Это нужно для обеспечения общественных фондов и условий для дальнейшего развития. Но общее направление движения общества к коммунизму, включающее и соответствующее развитие производительных сил, предполагает сведение такого труда ко всё меньшим и меньшим значениям. За известными пределами такой труд можно вообще не брать в расчёт, так как он уже не оказывает сколько-нибудь серьёзного влияния на общий ход человеческой жизни. Самое важное в необходимом труде при переходе к коммунизму, о котором здесь шла речь то, что этот, ещё не до конца свободный труд, является трудом по производству свободного времени и человека, способного действовать свободно в общественном масштабе.
\end{document}
